
\documentclass{article}
\usepackage{colortbl}
\usepackage{makecell}
\usepackage{multirow}
\usepackage{supertabular}

\begin{document}

\newcounter{utterance}

\centering \large Interaction Transcript for game `wordle', experiment `medium\_frequency\_words\_no\_clue\_no\_critic', episode 6 with reka{-}flash{-}3.1{-}t0.0.
\vspace{24pt}

{ \footnotesize  \setcounter{utterance}{1}
\setlength{\tabcolsep}{0pt}
\begin{supertabular}{c@{$\;$}|p{.15\linewidth}@{}p{.15\linewidth}p{.15\linewidth}p{.15\linewidth}p{.15\linewidth}p{.15\linewidth}}
    \# & \multicolumn{2}{c}{Player} && \multicolumn{2}{c}{Game Master} \\
    \hline

    \theutterance \stepcounter{utterance}  
    & & & \multicolumn{4}{p{0.6\linewidth}}{
        \cellcolor[rgb]{0.9,0.9,0.9}{
            \makecell[{{p{\linewidth}}}]{
                \texttt{\tiny{[P1$\langle$GM]}}
                \texttt{You are a language wizard who likes to guess words by using the given rules.} \\
\\ 
\texttt{Welcome to Wordle! You have six attempts to guess the target word, a valid English word of five lowercase letters (a{-}z). Please use the tags "explanation:" and "guess:" to provide a concise explanation for each guess.} \\
\\ 
\texttt{For instance, if your guess is "apple", your response should be} \\
\texttt{explanation: this is a common five{-}letter English word, and I am starting my guess with this word.} \\
\texttt{guess: apple} \\
\\ 
\texttt{After each guess, your answer will be validated, and you will receive feedback indicating which letters are correct (green), which letters are correct but in the wrong position (yellow), and which letters are incorrect (red). This feedback can be useful in determining which letters to include or exclude in your next guess.} \\
\\ 
\texttt{For example, the feedback for "apple" might be:} \\
\texttt{guess\_feedback: a<yellow> p<yellow> p<green> l<yellow> e<red>} \\
\\ 
\texttt{The explanation should contain details about how the guess\_feedback is used to arrive at a new guess.} \\
\\ 
\texttt{Do not generate any other text.} \\
\\ 
\texttt{Let's begin with your first guess.} \\
            }
        }
    }
     \\ \\

    \theutterance \stepcounter{utterance}  
    & \multicolumn{4}{p{0.6\linewidth}}{
        \cellcolor[rgb]{0.561,0.737,0.561}{
            \makecell[{{p{\linewidth}}}]{
                \texttt{\tiny{[P1$\rangle$GM]}}
                \texttt{<reasoning>} \\
\texttt{I need to start with my first guess in Wordle. Let's think about the best strategy. Usually, people recommend starting with common words that have a good mix of vowels and consonants to maximize information. Words like "crane" or "soare" are often suggested because they cover a range of letters.} \\
\\ 
\texttt{Wait, but maybe "raise" is a good start too. It has two vowels (a and e) and three consonants (r, s, z). The letters are spread out, which might help in getting feedback. Alternatively, "trace" is another option. Let me think about which letters are most common in English words. Letters like E, A, R, T, O, I, N, S, H, D are common. So maybe a word that includes some of these.} \\
\\ 
\texttt{Another approach is to pick a word that has unique letters to avoid ambiguity. For example, if I guess "apple", but if there are duplicate letters, the feedback might be confusing. So maybe avoid words with repeated letters in the first guess. "CRANE" has C, R, A, N, E – all unique. That could be good. Or "SLATE" – S, L, A, T, E. Also unique letters.} \\
\\ 
\texttt{Wait, but maybe "CRANE" is a bit less common. Let me check. Alternatively, "TRACE" is a valid word. Let's see. The letters in TRACE are T, R, A, C, E. All unique. That's a solid starting point. It includes vowels A and E, and consonants T, R, C. Maybe that's a good first guess. Alternatively, "STARE" – but that's also possible. Hmm.} \\
\\ 
\texttt{Alternatively, "CRANE" is a good starting word. Let me think. The letters are C, R, A, N, E. It's a valid word. Let's go with that. The reasoning is that it has a mix of vowels and consonants, all unique letters, which helps in getting clear feedback. So the first guess would be "crane".} \\
\\ 
\texttt{Wait, but maybe "raise" is better. Let's see. R, A, I, S, E. That includes I and S, which are common letters. But "raise" has two vowels (A and E) and three consonants. Maybe that's also a good choice. Alternatively, "soare"} \\
            }
        }
    }
    & & \\ \\

    \theutterance \stepcounter{utterance}  
    & & & \multicolumn{2}{p{0.3\linewidth}}{
        \cellcolor[rgb]{0.9,0.9,0.9}{
            \makecell[{{p{\linewidth}}}]{
                \texttt{\tiny{[GM$|$GM]}}
                \texttt{Error: The response should always start with the keyword 'explanation:'} \\
            }
        }
    }
    & & \\ \\

    \theutterance \stepcounter{utterance}  
    & & & \multicolumn{2}{p{0.3\linewidth}}{
        \cellcolor[rgb]{0.9,0.9,0.9}{
            \makecell[{{p{\linewidth}}}]{
                \texttt{\tiny{[GM$|$GM]}}
                \texttt{game\_result = ABORT} \\
            }
        }
    }
    & & \\ \\

\end{supertabular}
}

\end{document}
